\subsection{Dependent Systems in Three Variables} 
There are two basic ways to have a dependent system of three equations.
\begin{icpic}[x=\dimtd,y=\dimtd]
\begin{scope}[xshift=-.2\textwidth]
	\filldraw[pcb!50!pcc,plane]  (\project{-1}{0}{-1}) --
												(\project{1}{0}{-1}) --
												(\project{1}{0}{1}) --
												(\project{-1}{0}{1}) --
												cycle;
	\filldraw[pca,plane]  (\project{ .75}{-.75}{-1}) -- 
												(\project{ .75}{-.75}{ 1}) --
												(\project{-.75}{.75}{ 1}) --
												(\project{-.75}{.75}{-1}) --
												cycle;
	\draw[black!75,plane] (\project{0}{0}{-1}) -- (\project{0}{0}{1});
\end{scope}
\begin{scope}[xshift=.2\textwidth]
	\filldraw[black!75,plane] (\project{-1}{0}{-1}) --
												(\project{1}{0}{-1}) --
												(\project{1}{0}{1}) --
												(\project{-1}{0}{1}) --
												cycle;
\end{scope}														
\end{icpic}
Two or all three planes are \makeblank.
\begin{align*}
\begin{systemLa}
3x&+&1y&-&2z&=&6\\
3x&+&1y&-&2z&=&6\\
5x&-&2y&+&z&=&-1
\end{systemLa}
&&
\begin{systemLa}
7x&-&2y&+&5z&=&8\\
7x&-&2y&+&5z&=&8\\
7x&-&2y&+&5z&=&8
\end{systemLa}
\end{align*}
In standard form, two or all three equations are \makeblank.
\begin{icpic}[x=\dimtd,y=\dimtd]
\filldraw[pca,plane]  (\project{-1}{-1}{0}) --
                      (\project{ 1}{-1}{0}) --
										  (\project{ 1}{ 1}{0}) --
											(\project{-1}{ 1}{0}) -- 
											cycle;
\filldraw[pcb,plane] 	(\project{-.5}{-1}{ sqrt(.75)}) --
                      (\project{-.5}{ 1}{ sqrt(.75)}) --
                      (\project{ .5}{ 1}{-sqrt(.75)}) --
                      (\project{ .5}{-1}{-sqrt(.75)}) --
                      cycle;
\filldraw[pcc,plane] 	(\project{ .5}{-1}{ sqrt(.75)}) --
                      (\project{ .5}{ 1}{ sqrt(.75)}) --
                      (\project{-.5}{ 1}{-sqrt(.75)}) --
                      (\project{-.5}{-1}{-sqrt(.75)}) --
                      cycle;
\draw[black!75,plane] (\project{0}{1}{0}) -- (\project{0}{-1}{0});
\end{icpic}
Here, \makeblank[1.5in] planes meet at \makeblank. 
\[\begin{systemLa}
4x&-&y&-&3z&=&-3&\\
2x&+&y&-&z&=&1&\\
5x&+&4y&-&2z&=&5&
\end{systemLa}\]
Written in standard form, this system looks fine. If we try to solve it, we will end up describing the \makeblank[1in]. %What happens when we try to solve?
%\begin{lpic}{}{12}
%\end{lpic}
%When we get a system of \makeblanksmall equations in \makeblanksmall variables, the new system is \makeblank[1.5in].